% --- Archon physics formalization source begin ---
% archon:physics
% archon:covers PhyXMiniProblems/problem_phyx_mini_0954.lean
% archon:source-report reports/phyx_mini/problem_phyx_mini_0954.source.json

\chapter{Physics problem phyx\_mini\_0954}
\label{ch:PhyXMiniProblems_problem_phyx_mini_0954}

\paragraph{Problem source.}
\#\# Physical scenario

A magnetic lift uses a cylinder with radius \textbackslash{}( R \textbackslash{}) and width \textbackslash{}( W \textbackslash{}) wrapped by conducting wire \textbackslash{}( N \textbackslash{}) times on a diameter. A uniform external magnetic field \textbackslash{}( \textbackslash{}vec\{B\} \textbackslash{}) points downward while current \textbackslash{}( I \textbackslash{}) flows in the sense indicated. A mass \textbackslash{}( M \textbackslash{}) is hung from the cylinder by a cable attached to its rim on the axis of the coil. The height above the lowest possible position of the mass is \textbackslash{}( h \textbackslash{}).

\#\# Question

For what values of \textbackslash{}( \textbackslash{}sigma \textbackslash{}) will the cylinder rotate more than \textbackslash{}( 180\textasciicircum{}\textbackslash{}circ \textbackslash{}) if the mass is released from rest at \textbackslash{}( h = h\_\{\textbackslash{}text\{top\}\} \textbackslash{})?

\#\# Answer choices

- A: A: \textbackslash{}sigma < \textbackslash{}frac\{1\}\{2\}\textbackslash{}left(1 + \textbackslash{}frac\{\textbackslash{}pi\}\{2\}
ight)

- B: B: \textbackslash{}sigma > \textbackslash{}frac\{1\}\{2\}\textbackslash{}left(1 + \textbackslash{}frac\{\textbackslash{}pi\}\{2\}
ight)

- C: C: \textbackslash{}sigma > \textbackslash{}frac\{1\}\{2\}\textbackslash{}left(1 + \textbackslash{}pi
ight)

- D: D: \textbackslash{}sigma = \textbackslash{}frac\{1\}\{2\}\textbackslash{}left(1 + \textbackslash{}frac\{\textbackslash{}pi\}\{2\}
ight)

\#\# Figure caption (auxiliary; use the image as primary evidence)

The image illustrates a physical setup involving a pulley system in a magnetic field. Here's a detailed description:

- **Pulley System**: The central component is a large, grey circular pulley. It has a radius labeled "R" and a vertical axis of rotation. There is a wire wrapped around it, and an upward current "I" is indicated on the wire.

- **Magnetic Field**: A blue arrow labeled "\textbackslash{}(\textbackslash{}vec\{B\}\textbackslash{})" points downward, indicating the direction of a magnetic field.

- **Blocks and Weights**: 
  - To the right of the pulley, there is an orange block labeled "M," which is connected by a line to the pulley. This represents a mass that is hanging and applying a force through gravity.
  - Below the pulley, there is another beige block connected by a line, which represents another mass or weight. 

- **Forces and Measurements**:
  - "W" is labeled on both sides of the pulley, possibly indicating width or force vectors (e.g., tension).
  - An "h" is labeled vertically from the ground to the orange block, indicating height.
  
These components collectively illustrate a system where a magnetic field interacts with a current-carrying wire on a pulley, affecting the motion of masses.

\#\# Dataset metadata

Electromagnetic Induction; Physical Model Grounding Reasoning, Multi-Formula Reasoning

\paragraph{Recorded answer/context.}
B: B: \textbackslash{}sigma > \textbackslash{}frac\{1\}\{2\}\textbackslash{}left(1 + \textbackslash{}frac\{\textbackslash{}pi\}\{2\}
ight)

\paragraph{Figure/image path.}
/root/proposal\_for\_physic/science-mango/phyx\_mini\_run/phyx\_data/test\_image/954.png

\paragraph{Formalization target.}
create a compiling Lean file with sorry bodies at `PhyXMiniProblems/problem\_phyx\_mini\_0954.lean`. The Lean declarations must preserve the physical quantities, dimensions or dimensional roles, figure labels, governing-law hypotheses, and final relation expressed by this problem.
Use Mathlib/Physlib names found through LeanExplore where available. If a physics API is missing, introduce faithful local abstractions rather than scalar placeholder aliases.

\begin{theorem}[Physics formalization target]
\label{thm:physics:phyx_mini_0954:target}
\lean{PhyXMiniProblems.ProblemPhyXMini0954.rotatesMoreThanHalfTurn_iff_sigma_gt_threshold}
\uses{def:physics:phyx-mini-0954:phyxminiproblems-problemphyxmini0954-nonnegativereadout, def:physics:phyx-mini-0954:phyxminiproblems-problemphyxmini0954-energyinjoules, def:physics:phyx-mini-0954:phyxminiproblems-problemphyxmini0954-magneticliftsetup, def:physics:phyx-mini-0954:phyxminiproblems-problemphyxmini0954-initialattachmentangle, def:physics:phyx-mini-0954:phyxminiproblems-problemphyxmini0954-totalpotentialenergyinjoules, def:physics:phyx-mini-0954:phyxminiproblems-problemphyxmini0954-mechanicalenergyinjoulesattime, def:physics:phyx-mini-0954:phyxminiproblems-problemphyxmini0954-orientedangulartravelattime, def:physics:phyx-mini-0954:phyxminiproblems-problemphyxmini0954-rotatesmorethanhalfturn, def:physics:phyx-mini-0954:phyxminiproblems-problemphyxmini0954-matchesproblemandprimaryfigure, def:physics:phyx-mini-0954:phyxminiproblems-problemphyxmini0954-hasphysicalmagneticliftparameters, def:physics:phyx-mini-0954:phyxminiproblems-problemphyxmini0954-satisfiescurrentloopanduniformfieldlaws, def:physics:phyx-mini-0954:phyxminiproblems-problemphyxmini0954-satisfiesrigidcylinderandcablegeometry, def:physics:phyx-mini-0954:phyxminiproblems-problemphyxmini0954-satisfiesmagneticandgravitationalenergylaws, def:physics:phyx-mini-0954:phyxminiproblems-problemphyxmini0954-satisfiesconservativerotationaldynamics}
The undefined source parameter \(\sigma\) does not support a numerical
threshold theorem.  The strongest grounded conclusion is instead the
explicit apparatus potential, equality of release energy with initial
potential energy, and the necessary energy condition for every angular travel
between zero and a half-turn whenever the cylinder travels beyond a
half-turn.
\end{theorem}
\begin{proof}
Choose the fixed offset between attachment angle and magnetic-moment angle
from the rigid-geometry law.  Substitute successively the rim-height formula
\(h(\theta)=R(1-\cos\theta)\), the gravitational and magnetic potential laws,
the loop-area identity \(A=2RW\), the magnetic-moment identity
\(\mu=NIA\), and the fixed offset.  This yields
\[
 U(\theta)=MgR(1-\cos\theta)
   -(NI\,2RW)B\cos(\theta+\theta_0).
\]
At release the angular velocity is zero.  The kinetic-energy law therefore
makes the initial kinetic energy zero, so the initial mechanical energy is
the potential at the initial attachment angle.

Now assume the cylinder travels farther than \(\pi\) in the selected release
sense.  Oriented travel is continuous, equals zero at release, and exceeds
\(\pi\) at the witnessing positive time.  The intermediate value theorem
therefore supplies, for every \(s\in[0,\pi]\), a nonnegative time at which the
oriented travel is exactly \(s\).  The release sign is \(1\) or \(-1\), so
the attachment angle there is the initial angle plus the release sign times
\(s\).  Conservation of mechanical energy and the preceding initial-energy
identity give
\[
 K=U(\theta_{\rm initial})
   -U(\theta_{\rm initial}+\operatorname{sign}s).
\]
Finally the kinetic-energy law writes \(K\) as a positive inertia times a
square divided by two, hence \(K\geq0\); rearranging gives the asserted
potential-energy inequality.
\end{proof}

% --- Archon declaration topology begin ---
\section{Declaration topology}
\begin{definition}[electric Current Dimension]
  \label{def:physics:phyx-mini-0954:phyxminiproblems-problemphyxmini0954-electriccurrentdimension}
  \lean{PhyXMiniProblems.ProblemPhyXMini0954.electricCurrentDimension}
  Electric current has dimension charge per unit time
\end{definition}
\begin{proof}
  This is the declared model definition of electric Current Dimension.
\end{proof}
\begin{definition}[magnetic Flux Density Dimension]
  \label{def:physics:phyx-mini-0954:phyxminiproblems-problemphyxmini0954-magneticfluxdensitydimension}
  \lean{PhyXMiniProblems.ProblemPhyXMini0954.magneticFluxDensityDimension}
  Magnetic flux density (tesla) has dimension M T⁻¹ C⁻¹
\end{definition}
\begin{proof}
  This is the declared model definition of magnetic Flux Density Dimension.
\end{proof}
\begin{definition}[acceleration Dimension]
  \label{def:physics:phyx-mini-0954:phyxminiproblems-problemphyxmini0954-accelerationdimension}
  \lean{PhyXMiniProblems.ProblemPhyXMini0954.accelerationDimension}
  Acceleration has dimension L T⁻²
\end{definition}
\begin{proof}
  This is the declared model definition of acceleration Dimension.
\end{proof}
\begin{definition}[magnetic Moment Dimension]
  \label{def:physics:phyx-mini-0954:phyxminiproblems-problemphyxmini0954-magneticmomentdimension}
  \lean{PhyXMiniProblems.ProblemPhyXMini0954.magneticMomentDimension}
  A current loop's magnetic dipole moment has dimension C L² T⁻¹
\end{definition}
\begin{proof}
  This is the declared model definition of magnetic Moment Dimension.
\end{proof}
\begin{definition}[moment Of Inertia Dimension]
  \label{def:physics:phyx-mini-0954:phyxminiproblems-problemphyxmini0954-momentofinertiadimension}
  \lean{PhyXMiniProblems.ProblemPhyXMini0954.momentOfInertiaDimension}
  Moment of inertia has dimension M L²
\end{definition}
\begin{proof}
  This is the declared model definition of moment Of Inertia Dimension.
\end{proof}
\begin{definition}[torque Dimension]
  \label{def:physics:phyx-mini-0954:phyxminiproblems-problemphyxmini0954-torquedimension}
  \lean{PhyXMiniProblems.ProblemPhyXMini0954.torqueDimension}
  Torque has the same dimension as energy, M L² T⁻²
\end{definition}
\begin{proof}
  This is the declared model definition of torque Dimension.
\end{proof}
\begin{definition}[angular Velocity Dimension]
  \label{def:physics:phyx-mini-0954:phyxminiproblems-problemphyxmini0954-angularvelocitydimension}
  \lean{PhyXMiniProblems.ProblemPhyXMini0954.angularVelocityDimension}
  Angular velocity is an inverse-time quantity; radians are dimensionless
\end{definition}
\begin{proof}
  This is the declared model definition of angular Velocity Dimension.
\end{proof}
\begin{definition}[Length Magnitude]
  \label{def:physics:phyx-mini-0954:phyxminiproblems-problemphyxmini0954-lengthmagnitude}
  \lean{PhyXMiniProblems.ProblemPhyXMini0954.LengthMagnitude}
  A nonnegative, unit-independent physical length
\end{definition}
\begin{proof}
  This is the declared model definition of Length Magnitude.
\end{proof}
\begin{definition}[Area Magnitude]
  \label{def:physics:phyx-mini-0954:phyxminiproblems-problemphyxmini0954-areamagnitude}
  \lean{PhyXMiniProblems.ProblemPhyXMini0954.AreaMagnitude}
  A nonnegative, unit-independent physical area
\end{definition}
\begin{proof}
  This is the declared model definition of Area Magnitude.
\end{proof}
\begin{definition}[Mass Magnitude]
  \label{def:physics:phyx-mini-0954:phyxminiproblems-problemphyxmini0954-massmagnitude}
  \lean{PhyXMiniProblems.ProblemPhyXMini0954.MassMagnitude}
  A nonnegative, unit-independent mass
\end{definition}
\begin{proof}
  This is the declared model definition of Mass Magnitude.
\end{proof}
\begin{definition}[Electric Current Magnitude]
  \label{def:physics:phyx-mini-0954:phyxminiproblems-problemphyxmini0954-electriccurrentmagnitude}
  \lean{PhyXMiniProblems.ProblemPhyXMini0954.ElectricCurrentMagnitude}
  \uses{def:physics:phyx-mini-0954:phyxminiproblems-problemphyxmini0954-electriccurrentdimension}
  A nonnegative, unit-independent electric-current magnitude
\end{definition}
\begin{proof}
  This is the declared model definition of Electric Current Magnitude.
\end{proof}
\begin{definition}[Magnetic Flux Density Magnitude]
  \label{def:physics:phyx-mini-0954:phyxminiproblems-problemphyxmini0954-magneticfluxdensitymagnitude}
  \lean{PhyXMiniProblems.ProblemPhyXMini0954.MagneticFluxDensityMagnitude}
  \uses{def:physics:phyx-mini-0954:phyxminiproblems-problemphyxmini0954-magneticfluxdensitydimension}
  A nonnegative, unit-independent magnetic-flux-density magnitude
\end{definition}
\begin{proof}
  This is the declared model definition of Magnetic Flux Density Magnitude.
\end{proof}
\begin{definition}[Acceleration Magnitude]
  \label{def:physics:phyx-mini-0954:phyxminiproblems-problemphyxmini0954-accelerationmagnitude}
  \lean{PhyXMiniProblems.ProblemPhyXMini0954.AccelerationMagnitude}
  \uses{def:physics:phyx-mini-0954:phyxminiproblems-problemphyxmini0954-accelerationdimension}
  A nonnegative, unit-independent acceleration magnitude
\end{definition}
\begin{proof}
  This is the declared model definition of Acceleration Magnitude.
\end{proof}
\begin{definition}[Magnetic Moment Magnitude]
  \label{def:physics:phyx-mini-0954:phyxminiproblems-problemphyxmini0954-magneticmomentmagnitude}
  \lean{PhyXMiniProblems.ProblemPhyXMini0954.MagneticMomentMagnitude}
  \uses{def:physics:phyx-mini-0954:phyxminiproblems-problemphyxmini0954-magneticmomentdimension}
  A nonnegative, unit-independent magnetic dipole moment magnitude
\end{definition}
\begin{proof}
  This is the declared model definition of Magnetic Moment Magnitude.
\end{proof}
\begin{definition}[Moment Of Inertia Magnitude]
  \label{def:physics:phyx-mini-0954:phyxminiproblems-problemphyxmini0954-momentofinertiamagnitude}
  \lean{PhyXMiniProblems.ProblemPhyXMini0954.MomentOfInertiaMagnitude}
  \uses{def:physics:phyx-mini-0954:phyxminiproblems-problemphyxmini0954-momentofinertiadimension}
  A nonnegative, unit-independent moment of inertia
\end{definition}
\begin{proof}
  This is the declared model definition of Moment Of Inertia Magnitude.
\end{proof}
\begin{definition}[Axial Torque Quantity]
  \label{def:physics:phyx-mini-0954:phyxminiproblems-problemphyxmini0954-axialtorquequantity}
  \lean{PhyXMiniProblems.ProblemPhyXMini0954.AxialTorqueQuantity}
  \uses{def:physics:phyx-mini-0954:phyxminiproblems-problemphyxmini0954-torquedimension}
  A signed axial torque
\end{definition}
\begin{proof}
  This is the declared model definition of Axial Torque Quantity.
\end{proof}
\begin{definition}[Angular Velocity Quantity]
  \label{def:physics:phyx-mini-0954:phyxminiproblems-problemphyxmini0954-angularvelocityquantity}
  \lean{PhyXMiniProblems.ProblemPhyXMini0954.AngularVelocityQuantity}
  \uses{def:physics:phyx-mini-0954:phyxminiproblems-problemphyxmini0954-angularvelocitydimension}
  A signed angular velocity
\end{definition}
\begin{proof}
  This is the declared model definition of Angular Velocity Quantity.
\end{proof}
\begin{definition}[Energy Quantity]
  \label{def:physics:phyx-mini-0954:phyxminiproblems-problemphyxmini0954-energyquantity}
  \lean{PhyXMiniProblems.ProblemPhyXMini0954.EnergyQuantity}
  A signed mechanical energy; this is Physlib's dimensioned energy type
\end{definition}
\begin{proof}
  This is the declared model definition of Energy Quantity.
\end{proof}
\begin{definition}[nonnegative Readout]
  \label{def:physics:phyx-mini-0954:phyxminiproblems-problemphyxmini0954-nonnegativereadout}
  \lean{PhyXMiniProblems.ProblemPhyXMini0954.nonnegativeReadout}
  Read a nonnegative dimensionful scalar in a coherent system of units
\end{definition}
\begin{proof}
  This is the declared model definition of nonnegative Readout.
\end{proof}
\begin{definition}[signed Readout]
  \label{def:physics:phyx-mini-0954:phyxminiproblems-problemphyxmini0954-signedreadout}
  \lean{PhyXMiniProblems.ProblemPhyXMini0954.signedReadout}
  Read a signed dimensionful scalar in a coherent system of units
\end{definition}
\begin{proof}
  This is the declared model definition of signed Readout.
\end{proof}
\begin{definition}[energy In Joules]
  \label{def:physics:phyx-mini-0954:phyxminiproblems-problemphyxmini0954-energyinjoules}
  \lean{PhyXMiniProblems.ProblemPhyXMini0954.energyInJoules}
  \uses{def:physics:phyx-mini-0954:phyxminiproblems-problemphyxmini0954-energyquantity, def:physics:phyx-mini-0954:phyxminiproblems-problemphyxmini0954-signedreadout}
  Coherent-SI readout of an energy, in joules
\end{definition}
\begin{proof}
  This is the declared model definition of energy In Joules.
\end{proof}
\begin{definition}[Figure Object]
  \label{def:physics:phyx-mini-0954:phyxminiproblems-problemphyxmini0954-figureobject}
  \lean{PhyXMiniProblems.ProblemPhyXMini0954.FigureObject}
  Physical objects visible in image 954.png
\end{definition}
\begin{proof}
  This is the declared model definition of Figure Object.
\end{proof}
\begin{definition}[Figure Label]
  \label{def:physics:phyx-mini-0954:phyxminiproblems-problemphyxmini0954-figurelabel}
  \lean{PhyXMiniProblems.ProblemPhyXMini0954.FigureLabel}
  Literal symbolic labels present in the source or the supplied raster
\end{definition}
\begin{proof}
  This is the declared model definition of Figure Label.
\end{proof}
\begin{definition}[Vertical Direction]
  \label{def:physics:phyx-mini-0954:phyxminiproblems-problemphyxmini0954-verticaldirection}
  \lean{PhyXMiniProblems.ProblemPhyXMini0954.VerticalDirection}
  Named vertical directions needed to transcribe the arrows in the figure
\end{definition}
\begin{proof}
  This is the declared model definition of Vertical Direction.
\end{proof}
\begin{definition}[Winding Geometry]
  \label{def:physics:phyx-mini-0954:phyxminiproblems-problemphyxmini0954-windinggeometry}
  \lean{PhyXMiniProblems.ProblemPhyXMini0954.WindingGeometry}
  The winding geometry described by the prose and shown on the cylinder
\end{definition}
\begin{proof}
  This is the declared model definition of Winding Geometry.
\end{proof}
\begin{definition}[Cable Attachment]
  \label{def:physics:phyx-mini-0954:phyxminiproblems-problemphyxmini0954-cableattachment}
  \lean{PhyXMiniProblems.ProblemPhyXMini0954.CableAttachment}
  The cable attachment described in the problem statement
\end{definition}
\begin{proof}
  This is the declared model definition of Cable Attachment.
\end{proof}
\begin{definition}[Depicted Current Sense]
  \label{def:physics:phyx-mini-0954:phyxminiproblems-problemphyxmini0954-depictedcurrentsense}
  \lean{PhyXMiniProblems.ProblemPhyXMini0954.DepictedCurrentSense}
  The current-arrow sense directly visible on the front winding segment
\end{definition}
\begin{proof}
  This is the declared model definition of Depicted Current Sense.
\end{proof}
\begin{definition}[Magnetic Lift Figure]
  \label{def:physics:phyx-mini-0954:phyxminiproblems-problemphyxmini0954-magneticliftfigure}
  \lean{PhyXMiniProblems.ProblemPhyXMini0954.MagneticLiftFigure}
  \uses{def:physics:phyx-mini-0954:phyxminiproblems-problemphyxmini0954-figureobject, def:physics:phyx-mini-0954:phyxminiproblems-problemphyxmini0954-figurelabel, def:physics:phyx-mini-0954:phyxminiproblems-problemphyxmini0954-verticaldirection, def:physics:phyx-mini-0954:phyxminiproblems-problemphyxmini0954-windinggeometry, def:physics:phyx-mini-0954:phyxminiproblems-problemphyxmini0954-cableattachment, def:physics:phyx-mini-0954:phyxminiproblems-problemphyxmini0954-depictedcurrentsense}
  Typed qualitative data copied from the primary raster. The pale lower mass is recorded as a ghost of the same mass at its lowest possible position, not as a second physical load
\end{definition}
\begin{proof}
  This is the declared model definition of Magnetic Lift Figure.
\end{proof}
\begin{definition}[Spatial Vector]
  \label{def:physics:phyx-mini-0954:phyxminiproblems-problemphyxmini0954-spatialvector}
  \lean{PhyXMiniProblems.ProblemPhyXMini0954.SpatialVector}
  Three-dimensional Euclidean vectors used for magnetic-field readouts
\end{definition}
\begin{proof}
  This is the declared model definition of Spatial Vector.
\end{proof}
\begin{definition}[downward Unit Vector]
  \label{def:physics:phyx-mini-0954:phyxminiproblems-problemphyxmini0954-downwardunitvector}
  \lean{PhyXMiniProblems.ProblemPhyXMini0954.downwardUnitVector}
  \uses{def:physics:phyx-mini-0954:phyxminiproblems-problemphyxmini0954-spatialvector}
  Unit coordinate vector chosen to represent the downward vertical
\end{definition}
\begin{proof}
  This is the declared model definition of downward Unit Vector.
\end{proof}
\begin{definition}[Magnetic Lift Setup]
  \label{def:physics:phyx-mini-0954:phyxminiproblems-problemphyxmini0954-magneticliftsetup}
  \lean{PhyXMiniProblems.ProblemPhyXMini0954.MagneticLiftSetup}
  \uses{def:physics:phyx-mini-0954:phyxminiproblems-problemphyxmini0954-lengthmagnitude, def:physics:phyx-mini-0954:phyxminiproblems-problemphyxmini0954-areamagnitude, def:physics:phyx-mini-0954:phyxminiproblems-problemphyxmini0954-massmagnitude, def:physics:phyx-mini-0954:phyxminiproblems-problemphyxmini0954-electriccurrentmagnitude, def:physics:phyx-mini-0954:phyxminiproblems-problemphyxmini0954-magneticfluxdensitymagnitude, def:physics:phyx-mini-0954:phyxminiproblems-problemphyxmini0954-accelerationmagnitude, def:physics:phyx-mini-0954:phyxminiproblems-problemphyxmini0954-magneticmomentmagnitude, def:physics:phyx-mini-0954:phyxminiproblems-problemphyxmini0954-momentofinertiamagnitude, def:physics:phyx-mini-0954:phyxminiproblems-problemphyxmini0954-axialtorquequantity, def:physics:phyx-mini-0954:phyxminiproblems-problemphyxmini0954-angularvelocityquantity, def:physics:phyx-mini-0954:phyxminiproblems-problemphyxmini0954-energyquantity, def:physics:phyx-mini-0954:phyxminiproblems-problemphyxmini0954-magneticliftfigure}
  The physical apparatus, its one-degree-of-freedom trajectory, and independent energy and torque observables. The source's undefined sigma parameter is deliberately absent because no stated law connects it to the apparatus
\end{definition}
\begin{proof}
  This is the declared model definition of Magnetic Lift Setup.
\end{proof}
\begin{definition}[initial Attachment Angle]
  \label{def:physics:phyx-mini-0954:phyxminiproblems-problemphyxmini0954-initialattachmentangle}
  \lean{PhyXMiniProblems.ProblemPhyXMini0954.initialAttachmentAngle}
  \uses{def:physics:phyx-mini-0954:phyxminiproblems-problemphyxmini0954-magneticliftsetup}
  Initial attachment angle, at the release time t = 0
\end{definition}
\begin{proof}
  This is the declared model definition of initial Attachment Angle.
\end{proof}
\begin{definition}[total Potential Energy In Joules]
  \label{def:physics:phyx-mini-0954:phyxminiproblems-problemphyxmini0954-totalpotentialenergyinjoules}
  \lean{PhyXMiniProblems.ProblemPhyXMini0954.totalPotentialEnergyInJoules}
  \uses{def:physics:phyx-mini-0954:phyxminiproblems-problemphyxmini0954-energyinjoules, def:physics:phyx-mini-0954:phyxminiproblems-problemphyxmini0954-magneticliftsetup}
  Mechanical potential energy at a specified attachment angle
\end{definition}
\begin{proof}
  This is the declared model definition of total Potential Energy In Joules.
\end{proof}
\begin{definition}[mechanical Energy In Joules At Time]
  \label{def:physics:phyx-mini-0954:phyxminiproblems-problemphyxmini0954-mechanicalenergyinjoulesattime}
  \lean{PhyXMiniProblems.ProblemPhyXMini0954.mechanicalEnergyInJoulesAtTime}
  \uses{def:physics:phyx-mini-0954:phyxminiproblems-problemphyxmini0954-energyinjoules, def:physics:phyx-mini-0954:phyxminiproblems-problemphyxmini0954-magneticliftsetup, def:physics:phyx-mini-0954:phyxminiproblems-problemphyxmini0954-totalpotentialenergyinjoules}
  Total mechanical energy at a specified time after release
\end{definition}
\begin{proof}
  This is the declared model definition of mechanical Energy In Joules At Time.
\end{proof}
\begin{definition}[oriented Angular Travel At Time]
  \label{def:physics:phyx-mini-0954:phyxminiproblems-problemphyxmini0954-orientedangulartravelattime}
  \lean{PhyXMiniProblems.ProblemPhyXMini0954.orientedAngularTravelAtTime}
  \uses{def:physics:phyx-mini-0954:phyxminiproblems-problemphyxmini0954-magneticliftsetup, def:physics:phyx-mini-0954:phyxminiproblems-problemphyxmini0954-initialattachmentangle}
  Signed angular travel measured in the release rotation sense
\end{definition}
\begin{proof}
  This is the declared model definition of oriented Angular Travel At Time.
\end{proof}
\begin{definition}[Rotates More Than Half Turn]
  \label{def:physics:phyx-mini-0954:phyxminiproblems-problemphyxmini0954-rotatesmorethanhalfturn}
  \lean{PhyXMiniProblems.ProblemPhyXMini0954.RotatesMoreThanHalfTurn}
  \uses{def:physics:phyx-mini-0954:phyxminiproblems-problemphyxmini0954-magneticliftsetup, def:physics:phyx-mini-0954:phyxminiproblems-problemphyxmini0954-orientedangulartravelattime}
  The requested motion event: at some positive time the cylinder's angular travel in the release sense is strictly more than one half-turn
\end{definition}
\begin{proof}
  This is the declared model definition of Rotates More Than Half Turn.
\end{proof}
\begin{definition}[Matches Problem And Primary Figure]
  \label{def:physics:phyx-mini-0954:phyxminiproblems-problemphyxmini0954-matchesproblemandprimaryfigure}
  \lean{PhyXMiniProblems.ProblemPhyXMini0954.MatchesProblemAndPrimaryFigure}
  \uses{def:physics:phyx-mini-0954:phyxminiproblems-problemphyxmini0954-nonnegativereadout, def:physics:phyx-mini-0954:phyxminiproblems-problemphyxmini0954-signedreadout, def:physics:phyx-mini-0954:phyxminiproblems-problemphyxmini0954-magneticliftsetup, def:physics:phyx-mini-0954:phyxminiproblems-problemphyxmini0954-initialattachmentangle}
  Prose and raster evidence. No threshold inequality, rotation event, or normalized-energy identity occurs in this structure
\end{definition}
\begin{proof}
  This is the declared model definition of Matches Problem And Primary Figure.
\end{proof}
\begin{definition}[Has Physical Magnetic Lift Parameters]
  \label{def:physics:phyx-mini-0954:phyxminiproblems-problemphyxmini0954-hasphysicalmagneticliftparameters}
  \lean{PhyXMiniProblems.ProblemPhyXMini0954.HasPhysicalMagneticLiftParameters}
  \uses{def:physics:phyx-mini-0954:phyxminiproblems-problemphyxmini0954-nonnegativereadout, def:physics:phyx-mini-0954:phyxminiproblems-problemphyxmini0954-magneticliftsetup}
  Positivity and nondegeneracy conditions for the magnetic lift
\end{definition}
\begin{proof}
  This is the declared model definition of Has Physical Magnetic Lift Parameters.
\end{proof}
\begin{definition}[Satisfies Current Loop And Uniform Field Laws]
  \label{def:physics:phyx-mini-0954:phyxminiproblems-problemphyxmini0954-satisfiescurrentloopanduniformfieldlaws}
  \lean{PhyXMiniProblems.ProblemPhyXMini0954.SatisfiesCurrentLoopAndUniformFieldLaws}
  \uses{def:physics:phyx-mini-0954:phyxminiproblems-problemphyxmini0954-nonnegativereadout, def:physics:phyx-mini-0954:phyxminiproblems-problemphyxmini0954-downwardunitvector, def:physics:phyx-mini-0954:phyxminiproblems-problemphyxmini0954-magneticliftsetup}
  The current loop has area 2 R W, moment magnitude N I A, and is immersed in a uniform downward field. These are general physical laws and do not mention the requested sigma threshold
\end{definition}
\begin{proof}
  This is the declared model definition of Satisfies Current Loop And Uniform Field Laws.
\end{proof}
\begin{definition}[Satisfies Rigid Cylinder And Cable Geometry]
  \label{def:physics:phyx-mini-0954:phyxminiproblems-problemphyxmini0954-satisfiesrigidcylinderandcablegeometry}
  \lean{PhyXMiniProblems.ProblemPhyXMini0954.SatisfiesRigidCylinderAndCableGeometry}
  \uses{def:physics:phyx-mini-0954:phyxminiproblems-problemphyxmini0954-nonnegativereadout, def:physics:phyx-mini-0954:phyxminiproblems-problemphyxmini0954-magneticliftsetup, def:physics:phyx-mini-0954:phyxminiproblems-problemphyxmini0954-orientedangulartravelattime}
  Geometry of the material attachment point fixed to the rim: at attachment angle theta measured from downward, the height is R(1-cos theta), the top height is 2R, and the coil axis has a fixed angular offset from the attachment coordinate
\end{definition}
\begin{proof}
  This is the declared model definition of Satisfies Rigid Cylinder And Cable Geometry.
\end{proof}
\begin{definition}[Satisfies Magnetic And Gravitational Energy Laws]
  \label{def:physics:phyx-mini-0954:phyxminiproblems-problemphyxmini0954-satisfiesmagneticandgravitationalenergylaws}
  \lean{PhyXMiniProblems.ProblemPhyXMini0954.SatisfiesMagneticAndGravitationalEnergyLaws}
  \uses{def:physics:phyx-mini-0954:phyxminiproblems-problemphyxmini0954-nonnegativereadout, def:physics:phyx-mini-0954:phyxminiproblems-problemphyxmini0954-signedreadout, def:physics:phyx-mini-0954:phyxminiproblems-problemphyxmini0954-magneticliftsetup}
  Potential-energy and torque laws for gravity and for a magnetic dipole in a uniform field. For the fixed rim attachment the gravitational torque is -M g R sin theta when increasing attachment angle is positive. The net torque is the sum of its two physical contributions
\end{definition}
\begin{proof}
  This is the declared model definition of Satisfies Magnetic And Gravitational Energy Laws.
\end{proof}
\begin{definition}[Satisfies Conservative Rotational Dynamics]
  \label{def:physics:phyx-mini-0954:phyxminiproblems-problemphyxmini0954-satisfiesconservativerotationaldynamics}
  \lean{PhyXMiniProblems.ProblemPhyXMini0954.SatisfiesConservativeRotationalDynamics}
  \uses{def:physics:phyx-mini-0954:phyxminiproblems-problemphyxmini0954-nonnegativereadout, def:physics:phyx-mini-0954:phyxminiproblems-problemphyxmini0954-signedreadout, def:physics:phyx-mini-0954:phyxminiproblems-problemphyxmini0954-magneticliftsetup, def:physics:phyx-mini-0954:phyxminiproblems-problemphyxmini0954-initialattachmentangle, def:physics:phyx-mini-0954:phyxminiproblems-problemphyxmini0954-totalpotentialenergyinjoules, def:physics:phyx-mini-0954:phyxminiproblems-problemphyxmini0954-mechanicalenergyinjoulesattime, def:physics:phyx-mini-0954:phyxminiproblems-problemphyxmini0954-rotatesmorethanhalfturn}
  Conservative one-dimensional rotational mechanics: continuity of angular position, the angular-velocity derivative relation, nonnegative quadratic kinetic energy, and conservation of total mechanical energy after release. It contains no reachability criterion, sigma condition, or answer-choice expression
\end{definition}
\begin{proof}
  This is the declared model definition of Satisfies Conservative Rotational Dynamics.
\end{proof}
\begin{definition}[Answer Choice]
  \label{def:physics:phyx-mini-0954:phyxminiproblems-problemphyxmini0954-answerchoice}
  \lean{PhyXMiniProblems.ProblemPhyXMini0954.AnswerChoice}
  Labels of the four displayed answer choices
\end{definition}
\begin{proof}
  This is the declared model definition of Answer Choice.
\end{proof}
\begin{definition}[half Turn Sigma Threshold]
  \label{def:physics:phyx-mini-0954:phyxminiproblems-problemphyxmini0954-halfturnsigmathreshold}
  \lean{PhyXMiniProblems.ProblemPhyXMini0954.halfTurnSigmaThreshold}
  The boundary printed in choices A, B, and D
\end{definition}
\begin{proof}
  This is the declared model definition of half Turn Sigma Threshold.
\end{proof}
\begin{definition}[displayed Sigma Condition]
  \label{def:physics:phyx-mini-0954:phyxminiproblems-problemphyxmini0954-displayedsigmacondition}
  \lean{PhyXMiniProblems.ProblemPhyXMini0954.displayedSigmaCondition}
  \uses{def:physics:phyx-mini-0954:phyxminiproblems-problemphyxmini0954-answerchoice, def:physics:phyx-mini-0954:phyxminiproblems-problemphyxmini0954-halfturnsigmathreshold}
  The four propositions literally displayed beside the answer labels
\end{definition}
\begin{proof}
  This is the declared model definition of displayed Sigma Condition.
\end{proof}
\begin{definition}[Recorded dataset answer]
  \label{def:physics:phyx-mini-0954:phyxminiproblems-problemphyxmini0954-recordeddatasetanswer}
  \lean{PhyXMiniProblems.ProblemPhyXMini0954.recordedDatasetAnswer}
  \uses{def:physics:phyx-mini-0954:phyxminiproblems-problemphyxmini0954-answerchoice}
  The private answer-sheet metadata records choice B, independently of the magnetic-lift dynamics.
\end{definition}
\begin{proof}
  This is the literal private metadata assignment; the block is retained because the Lean-aware graph scans this declaration.
\end{proof}
% --- Archon declaration topology end ---
% --- Archon physics formalization source end ---
